
\textbf{Proposition 1.7.} (\uline{Comparer} EGA IV 19.1.5 (iv)). \uline{Soient} @@M0@@, \uline{et} @@M1@@ \uline{deux immersions}:
\begin{enumerate}[label=\roman*)]
\item \uline{Si} @@M2@@ \uline{et} @@M3@@ \uline{sont régulières}, @@M4@@ \uline{est régulière}, \uline{et la suite}:
@@M5@@
\uline{est exacte}; \uline{en particulier}, @@M6@@.
\item \uline{Si} @@M7@@ \uline{et} @@M8@@ \uline{sont régulières}, \uline{et si la suite}:
@@M9@@
\uline{est exacte et localement scindée}, @@M10@@ \uline{est régulière};
\item \uline{Si} @@M11@@ \uline{est noethérien}, \uline{et si} @@M12@@ \uline{et} @@M13@@ \uline{sont régulières}, @@M14@@ \uline{est régulière aux points de} @@M15@@.
\end{enumerate}

La question étant locale sur @@M16@@, on peut supposer @@M17@@ affine et les immersions fermées; soient @@M18@@, @@M19@@, @@M20@@, avec @@M21@@, @@M22@@.

i) Quitte à se localiser, on peut supposer @@M23@@ libre sur @@M24@@, et @@M25@@ libre sur @@M26@@ (1.3). Soient @@M27@@ des éléments de @@M28@@ relevant une base de @@M29@@, @@M30@@ des éléments de @@M31@@ relevant une base de @@M32@@; alors, quitte à se localiser davantage, on peut supposer que @@M33@@ engendrent @@M34@@, que @@M35@@ engendrent @@M36@@, et que @@M37@@ engendrent @@M38@@, @@M39@@ étant l'image de @@M40@@ dans @@M41@@ (1.4); de plus, @@M42@@ forment un système minimum de générateurs de @@M43@@, et @@M44@@ un système minimum de générateurs de @@M45@@.

On considère le diagramme:
@@M46@@
